In the non-additive ramp loss model, we wish to replace
\[
    \zeta_i + \sum_{k \in \K} \gamma_{ik} w_k \geq 2, \tquad i \in \P,
\]
by
\[
    \zeta_i + \umax{k \in \K} \; \gamma_{ik} w_k \geq 2, \tquad i \in \P.
\]
which is equivalent to
\begin{equation*}
    \begin{split}
        \zeta_i - 2 &\geq - \umax{k \in \K} \; \gamma_{ik} w_k\\
        &= \umin{k \in \K} \; -\gamma_{ik} w_k , \tquad i \in \P. \\
    \end{split}
\end{equation*}
Define
\[
    \eta_{ik} = -\gamma_{ik}, \tquad i \in \P, \; k \in \K.
\]
Then one formulation of the non-additive ramp-loss model is
\begin{subequations}
    \label{eq:NRM}
    \begin{align}
        \mathbf{\min} \quad & \sum_{i \in \cP} \zeta_i + \sum_{i \in \cN} \sum_{k \in \K} \gamma_{ik} w_k \label{eq:NRM-obj} \\
        \mathbf{\st} \quad & \zeta_i - 2 \geq \eta_{ik'} - \sum_{\substack{k \in \K, \\ \eta_{ik} < \eta_{ik'}}} (\eta_{ik'}-\eta_{ik}) w_k,
        & \forall i \in \cP, \; k' \in \K, \label{eq:NRM-misP} \\
        & \zeta_i - 2 \geq - \sum_{k \in \K} 2w_k, & \forall i \in \I, \label{eq:NRM-fix}\\
        & \sum_{k \in \K} c_k w_k \leq C, & \label{eq:NRM-complex} \\
        & w \in \{0, 1\}^{\K}, \zeta \in [0,2]^{\cP}. \label{eq:NRM-vars}
    \end{align}
\end{subequations}
The reduced cost of a rule is derived in \autoref{sec:reduced_costs_nonadditive_ramp_loss}.
However, this reduced cost cannot be easily be modeled inside a MIP\@.

Alternatively, we may consider the following formulation:
\begin{subequations}
    \label{eq:NRMa}
    \begin{align}
        \mathbf{\min} \quad & \sum_{i \in \cP} \zeta_i + \sum_{i \in \cN} \sum_{k \in \K} \gamma_{ik} w_k \label{eq:NRMa-obj} \\
        \mathbf{\st} \quad & \zeta_i - 2 \geq \sum_{k \in \K} u_{ik}\eta_{ik}, & \forall i \in \cP, \label{eq:NRMa-misP} \\
        & \sum_{k \in \K} u_{ik} = 1, & \forall i \in \cP, \label{eq:NRMa-misP2} \\
        & u_{ik} \leq w_k, & \forall i \in \cP, \; k \in \K, \label{eq:NRMa-misP3} \\
        & \sum_{k \in \K} c_k w_k \leq C, & \label{eq:NRMa-complex} \\
        & w \in \{0, 1\}^{\K}, \zeta \in [0,2]^{\cP}, u \in [0,1]^{\P \times \K}. \label{eq:NRMa-vars}
    \end{align}
\end{subequations}
While this formulation has a large number of variables, the $u$-variables need not be binary.
The reduced cost of a variable $w_k,\; k \in \K$, is derived in \autoref{sec:reduced_costs_nonadditive_ramp_loss}.
Implementing column-and-row generation for this master problem formulation is not trivial,
as a rule corresponds to a $w$-variable and many $u$-variables.

\begin{note}
    A formulation similar to \eqref{eq:NRMa} may be constructed for the pricing problem \eqref{eq:RP_set_formulation}.
    Implementing and solving this MIP for a number of test cases revealed that it is not easily solved by SCIP.\@
    Using proprietary solver Gurobi, as well as some tricks exploiting the nature of the data,
    performance improved moderately, but not enough to solve the problem in a handful of minutes.
    Hence, the more compact formulation \eqref{eq:RP} was explored in combination with separation of cutting planes.
    Ultimately, none of the approaches resulted in the ability to solve the pricing problem efficiently for our dataset.
\end{note}

Evidently, solving the relaxed non-additive ramp loss model exactly is difficult,
as even the LP cannot be easily solved via column (and row) generation.
A pragmatic approach is to rely solely on a pricing heuristic.
Adapting the heuristic pricer to the non-additive model is possible, but requires substantial work.
Therefore, we first explore whether the non-additive model gives better results on example instances.
We consider again the setting and examples in \autoref{sec:ramp-examples}.
We use formulation \eqref{eq:NRM}, as it is easiest to implement within the existing implementation of the ramp loss model.

\begin{note}
    In the construction and implementation of \eqref{eq:NRM}, constraint \eqref{eq:NRM-fix} was overlooked.
    This affects the objective of the empty ruleset if there exist $i \in \P$ such that $\eta_{ik} < 0$ for all
    $k$ considered in the branching phase.
    The objective of a non-empty ruleset is correctly modeled.
    Since we are interested in an empty ruleset, and non-empty rulesets are indeed found as solutions,
    the mistake of omitting \eqref{eq:NRM-fix} seems to have no effect in practice.
\end{note}

\begin{example}[Two clouds, non-additive model]\label{ex:two_clouds_nonadditive}
    Consider again \autoref{ex:two_clouds} and the tested parameters.

    For $C=1$ and any $\beta$, the results are identical to those of the additive model.
    For $C=2,$ $\beta=0,$ the ruleset differs and is depicted in \autoref{fig:two_clouds_nonadditive_solution_bad_a}.
    It consists of two rules, and the covered region is identical to the one-rule solution found for the additive model.
    The objective is also identical.
    To the model, both solutions are equally good.
    For classification tasks, they are also identical.
    If the smaller ruleset is preferred, for instance because of interpretability, then $C$ could be chosen as low as possible so long as the performance,
    whichever way it is measured, does not decrease.
    Another solution could be to penalize ruleset complexity in the model.
    It is important to note that this solution, when considered in the additive model with $\beta =0$,
    has the exact same objective value as \autoref{fig:two_clouds_solution_simple}.
    Hence, the choice of this solution in the non-additive model is a coincidence rather than an artifact.
    In both the additive and non-additive model, the value of $C$ is best kept small to prevent this.

    For $C=2,$ $\beta=0.5$ the solution is similar to \autoref{fig:two_clouds_solution_bad_a} and is shown in \autoref{fig:two_clouds_nonadditive_solution_bad_b}.
    For $C=2,$ $\beta=1$, the solution contains $x_1 \geq 9.88$, but also another rule, as shown in \autoref{fig:two_clouds_nonadditive_solution_bad_c}.
    Removing this other rule from the ruleset, leaving just the rule $x_1 \geq 9.88$, results in the same objective.
    Hence, one may want to prune solutions.
    One way of pruning solutions is to consider all subsets of the ruleset, and pick the smallest subset that has the same objective value as the full ruleset.
    Alternatively, one could again limit $C$ or penalize rule complexity.

    Notably, for $\beta > 0.5$, the non-additive model seems less included to produce undesirable rulesets than the additive model.
    While the additive model yields two-rule solutions for $C=2,$ $\beta>0.5,$
    the non-additive yields a one-rule solution either immediately or after pruning.
    This can be partly attributed to the non-additivity that was introduced.
    Another factor is the negative points' loss, which is calculated in the same way in either model.
    The loss for negative points is additive, hence using more rules is bad for negative points in or around the covered region.

    \begin{figure}
    \centering
    \begin{subfigure}{.48\textwidth}
        \centering
        \includesvg[width=\textwidth]{gen/test_clouds/results/O_1_C_2_mw_0_nonadditive.set.svg}
        \caption{}
        \label{fig:two_clouds_nonadditive_solution_bad_a}
    \end{subfigure}
    \hfill
    \begin{subfigure}{.48\textwidth}
        \centering
        \includesvg[width=\textwidth]{gen/test_clouds/results/O_1_C_2_mw_0.5_nonadditive.set.svg}
        \caption{}
        \label{fig:two_clouds_nonadditive_solution_bad_b}
    \end{subfigure} \\
        \begin{subfigure}{.48\textwidth}
        \centering
        \includesvg[width=\textwidth]{gen/test_clouds/results/O_1_C_2_mw_1_nonadditive.set.svg}
        \caption{}
        \label{fig:two_clouds_nonadditive_solution_bad_c}
    \end{subfigure} \\
    \caption{Solutions to an instance of \autoref{ex:two_clouds} computed with the ramp loss model, configured with $C=2$, $\beta=0.5$ (a) and $C=2$, $\beta=1$ (b), respectively.}
    \label{fig:two_clouds_nonadditive_solution_bad}
\end{figure}

\end{example}

We conclude that the non-additive ramp loss model behaves better than the additive ramp loss model.
In practice, we perform pricing on the additive variant.
Then, we use the non-additive variant inside the branching stage.
This effectively means we take the columns generated using the additive variant,
add these columns to an IP of the non-additive variant, and solve.